% ============================================================================
% SAFETY ARCHITECTURE
% ============================================================================

\section{Safety Architecture}
\label{section:safety-architecture}

The STAR robot has actuators that can move heavy chassis at meters
per second; "fail safe" is not optional. This section describes the
overlapping safety mechanisms so a future engineer can identify, at
review time, which layer actually owns a given safety guarantee.

\subsection{Defense in depth, in three layers}

\begin{enumerate}
  \item \textbf{ROS2 supervisor (Pi 5).} \texttt{star\_supervisor}
        gates command flow between manual / autonomous sources and
        the hardware bridge. It is the only thing that publishes on
        \texttt{/cmd\_vel\_out}.
  \item \textbf{Gateway and firmware comm layer (Pi 5 + RX72N).} The
        gateway forwards \texttt{/star/estop} into the protobuf
        \texttt{EmergencyStopRequest} message; the firmware
        \texttt{comm\_task} immediately writes
        \texttt{shared\_data.estop\_active = true}.
  \item \textbf{Firmware motor / watchdog layer (RX72N).}
        \texttt{motor\_control\_task} forces zero duty when
        \texttt{estop\_active} is set; the IWDT independently resets
        the chip if the watchdog task ever stops kicking; DRV8263
        \texttt{nFAULT} / OLP detection triggers a per-motor disable.
\end{enumerate}

A single bug in any one layer cannot, by itself, cause the robot to
move when the operator has called e-stop.

\subsection{The supervisor model}

The robot hardware bridge reads \texttt{/cmd\_vel\_out}, never
\texttt{/cmd\_vel} directly. \texttt{star\_supervisor} arbitrates:

\begin{center}
\begin{tabular}{|l|l|}
\hline
\textbf{Source topic}                      & \textbf{Forwarded to /cmd\_vel\_out when}                       \\
\hline
\texttt{/cmd\_vel} (manual)                & \texttt{autonomy\_enable=false} AND \texttt{estop=false}        \\
\texttt{/nav2/cmd\_vel}                    & \texttt{autonomy\_enable=true}  AND \texttt{estop=false}        \\
\texttt{/goal\_pose} -> \texttt{/nav2/goal\_pose} & \texttt{autonomy\_enable=true} AND \texttt{estop=false}    \\
\hline
\end{tabular}
\end{center}

\texttt{/star/autonomy\_enable} and \texttt{/star/estop} are both
latched booleans (\texttt{TransientLocal} QoS) so a node that joins
late still sees the current state.

When \texttt{estop=true} the supervisor publishes a zero
\texttt{Twist} on \texttt{/cmd\_vel\_out} continuously (not just
once), so the hardware bridge cannot drift back to a stale non-zero
command if a publisher misbehaves.

\subsection{Estop signal flow}

\begin{verbatim}
operator clicks "STOP" in UI / safety_monitor trips a sonar threshold
        |
        v
ROS2 publishes /star/estop = true (TransientLocal latched)
        |
        +---> star_supervisor: zeros /cmd_vel_out (Layer 1)
        |
        +---> star_gateway_bridge -> Go gateway:
        |       sends EmergencyStopRequest over USB CDC / SPI
        |
        v
RX72N comm_task receives, calls internal_handle_estop,
        sets shared_data.estop_active = true (Layer 2)
        |
        v
RX72N motor_control_task next tick:
        if (shared_data.estop_active) duty = 0; state = ESTOP (Layer 3)
        rx_motor_emergency_stop() drives DRVOFF high on every
        DRV8263 channel as a hardware-level disable
\end{verbatim}

Recovery requires explicit operator action: clear the latched
\texttt{/star/estop} (UI button, or \texttt{ros2 topic pub}). The
firmware does not auto-clear.

\subsection{Watchdog architecture}

The RX72N has two watchdog peripherals:

\begin{center}
\begin{tabular}{|l|p{5.0cm}|p{5.5cm}|}
\hline
\textbf{Feature}        & \textbf{IWDT (Independent)}                & \textbf{WDT}                                \\
\hline
Clock source            & 15\,kHz dedicated low-speed oscillator     & PCLKB (60\,MHz)                              \\
Stoppable in software   & No, always on                              & Yes                                          \\
Use case in STAR        & Hard hang detection                        & Reserved for soft-debug timeouts             \\
Reset flag              & RSTSR2.IWDTRF                              & RSTSR2.WDTRF                                 \\
On-reset behavior       & \texttt{internal\_rx\_fatal\_error} halts  & Logs warning, returns boot status code        \\
Period                  & ~1\,s with current divider                  & --                                            \\
\hline
\end{tabular}
\end{center}

\texttt{watchdog\_monitor\_task} runs at 200\,Hz and:

\begin{enumerate}
  \item Reads each task's last-tick timestamp from
        \texttt{shared\_data.task\_heartbeats}.
  \item If any task is more than 3 ticks late, logs an error and
        asserts \texttt{shared\_data.estop\_active}.
  \item Calls \texttt{rx\_iwdt\_feed()} only after every check has
        passed.
  \item Reads stack high-watermarks and asserts on starvation
        (less than 256\,B remaining on any task).
\end{enumerate}

Because the IWDT can never be stopped, a permanently-stuck
\texttt{watchdog\_monitor\_task} can only fail to kick, after which
the IWDT will reset the chip. The firmware on next boot reads
RSTSR2.IWDTRF, logs the failure, and halts so a debugger can read
the post-mortem state.

\subsection{Hardware-level safety: DRV8263 + DRVOFF}

Each DRV8263H motor driver has three control pins:

\begin{itemize}
  \item \texttt{nSLEEP}: low forces the chip into sleep, all outputs
        Hi-Z. Pulled high during normal operation.
  \item \texttt{DRVOFF}: high forces all outputs Hi-Z without
        sleeping the chip. Used as a fast disable on fault.
  \item \texttt{nFAULT}: open-drain output asserted by the chip on
        overcurrent / thermal / OLP / Vm undervoltage. RX72N reads it
        as a GPIO input.
\end{itemize}

\texttt{rx\_motor\_emergency\_stop} drives DRVOFF high before
attempting any PWM cleanup, so even if the GPTW PWM keeps running,
the motor is electrically isolated within microseconds.

The Open Load Protection (OLP) diagnostic in \texttt{rx\_drv8263}
runs at startup and on operator request; it cycles a 3-pattern test
across each motor with DRVOFF high to detect open / shorted loads
before they become runtime faults.

\subsection{ThreadX FPU and stack safety}

ThreadX RXv3 port supports per-thread FPU context, but only for
threads that opt in. Tasks that touch \texttt{float} / \texttt{double}
must declare \texttt{TX\_THREAD\_EXTENSION\_2} and be created with
\texttt{TX\_ENABLE\_FPU\_SUPPORT}. If a task uses floats without the
extension, you get silent corruption when another FPU-using task
preempts it.

Every motor / PID / IMU / telemetry task in STAR enables FPU
support; the LED status / serial bringup tasks do not. The
\texttt{watchdog\_monitor\_task} verifies stack high-watermarks on
every tick to catch overflow before it corrupts another task's stack.

\subsection{Lifecycle nodes (ROS2 side)}

\texttt{star\_safety\_monitor} is a \texttt{LifecycleNode}; it
transitions:

\begin{verbatim}
UNCONFIGURED  --on_configure-->  INACTIVE
INACTIVE      --on_activate-->   ACTIVE
ACTIVE        --on_deactivate--> INACTIVE
INACTIVE      --on_cleanup-->    UNCONFIGURED
\end{verbatim}

While \texttt{INACTIVE} the node listens for diagnostics but does
not publish \texttt{/star/estop}. While \texttt{ACTIVE} it monitors
sonar topics, IMU jolts, and the supervisor's bond, and trips the
estop when any threshold breaches. The \texttt{on\_shutdown}
callback gracefully publishes \texttt{/star/estop=true} as a final
"safe stop" before exit.

\textbf{Reconfigurability:} \texttt{on\_cleanup()} is required by REP-2003
to leave the node in a state where a subsequent \texttt{on\_configure()}
succeeds. The implementation \texttt{undeclare\_parameter()}s all 12
declared parameters in addition to releasing publishers, subscribers
and the bond -- otherwise the second \texttt{declare\_parameter()} call
would throw \texttt{AlreadyDeclaredParameterException}.

\textbf{Safe-rate guard:} \texttt{on\_activate()} clamps
\texttt{publish\_rate} to a $\geq 0.1\,\text{Hz}$ floor (10 s monitoring
period) and rejects \texttt{NaN}/$\pm\infty$. A user override of $0.0$
would otherwise produce a 0 ms timer (busy-spin) via the
$1000.0/\text{publish\_rate}$ division -- the safety monitor must
never be the cause of a CPU peg. The default rate is $10\,\text{Hz}$.

\subsection{Where to look in code}

\begin{center}
\begin{tabular}{|p{5.0cm}|p{9.0cm}|}
\hline
\textbf{Topic}                                  & \textbf{Source}                                                                                              \\
\hline
ROS2 supervisor / arbitration                    & \texttt{star-ros2/src/star\_supervisor/}                                                                     \\
ROS2 lifecycle safety monitor                    & \texttt{star-ros2/src/star\_safety\_monitor/}                                                                \\
Estop forwarding (gateway)                       & \texttt{star-gateway/internal/service/}                                                                       \\
Firmware estop hot path                          & \texttt{star-rx72n-firmware/src/tasks/motor\_control\_task.c}                                                 \\
Firmware estop receive path                      & \texttt{star-rx72n-firmware/src/tasks/comm\_task.c}                                                           \\
Watchdog monitor                                  & \texttt{star-rx72n-firmware/src/tasks/watchdog\_monitor\_task.c}                                              \\
IWDT driver                                       & \texttt{star-rx72n-firmware/libs/rx\_core/src/rx\_iwdt.c}                                                     \\
DRV8263 hardware-level disable                   & \texttt{star-rx72n-firmware/libs/rx\_drv8263/} + DRVOFF GPIO in hardware\_config.h                            \\
RSTSR2 reset-cause inspection                    & \texttt{star-rx72n-firmware/src/main.c}                                                                       \\
ARQ -> HARQ migration                            & \texttt{docs/decisions/ADR-001-arq-to-harq.md}                                                                \\
\hline
\end{tabular}
\end{center}
